\documentclass[11pt,a4paper]{article}
% preamble.tex — estilo común de todos los apuntes Froth.
% Cada apunte hace % preamble.tex — estilo común de todos los apuntes Froth.
% Cada apunte hace % preamble.tex — estilo común de todos los apuntes Froth.
% Cada apunte hace \input{preamble} justo después de \documentclass.
\usepackage[margin=2.4cm]{geometry}
\usepackage{xcolor}
\definecolor{litblue}{RGB}{31,79,121}
\definecolor{litgreen}{RGB}{27,94,32}
\definecolor{litorange}{RGB}{230,81,0}
\usepackage{parskip}
\usepackage{titlesec}
\titleformat{\section}{\Large\bfseries\color{litblue}}{\thesection}{0.6em}{}
\titleformat{\subsection}{\large\bfseries\color{litblue}}{\thesubsection}{0.6em}{}
\usepackage{tcolorbox}
\usepackage{fancyhdr}
\pagestyle{fancy}
\fancyhf{}
\fancyhead[L]{\small\color{litblue}Froth — Apuntes de ML}
\fancyhead[R]{\small\thepage}
\renewcommand{\headrulewidth}{0.4pt}
\usepackage[colorlinks=true,linkcolor=litblue,urlcolor=litblue]{hyperref}

% Cabecera de cada apunte: \cabecera{numero}{titulo}
\newcommand{\cabecera}[2]{%
  \begin{tcolorbox}[colback=litblue,coltext=white,colframe=litblue,arc=2mm]
    {\Large\bfseries Apunte #1 — #2}\\[3pt]
    {\small Proyecto Froth · aprender Machine Learning construyendo}
  \end{tcolorbox}\vspace{0.8em}}

% Cajas temáticas
\newtcolorbox{concepto}[1]{colback=blue!4,colframe=litblue,title={Concepto: #1},arc=1.5mm}
\newtcolorbox{practica}{colback=green!5,colframe=litgreen,title={Pruébalo tú (teclado en mano)},arc=1.5mm}
\newtcolorbox{ojo}{colback=orange!8,colframe=litorange,title={Ojo / error típico},arc=1.5mm}
\newtcolorbox{resumen}{colback=gray!8,colframe=gray!60,title={En una frase},arc=1.5mm}

\newcommand{\codigo}[1]{\texttt{#1}}
 justo después de \documentclass.
\usepackage[margin=2.4cm]{geometry}
\usepackage{xcolor}
\definecolor{litblue}{RGB}{31,79,121}
\definecolor{litgreen}{RGB}{27,94,32}
\definecolor{litorange}{RGB}{230,81,0}
\usepackage{parskip}
\usepackage{titlesec}
\titleformat{\section}{\Large\bfseries\color{litblue}}{\thesection}{0.6em}{}
\titleformat{\subsection}{\large\bfseries\color{litblue}}{\thesubsection}{0.6em}{}
\usepackage{tcolorbox}
\usepackage{fancyhdr}
\pagestyle{fancy}
\fancyhf{}
\fancyhead[L]{\small\color{litblue}Froth — Apuntes de ML}
\fancyhead[R]{\small\thepage}
\renewcommand{\headrulewidth}{0.4pt}
\usepackage[colorlinks=true,linkcolor=litblue,urlcolor=litblue]{hyperref}

% Cabecera de cada apunte: \cabecera{numero}{titulo}
\newcommand{\cabecera}[2]{%
  \begin{tcolorbox}[colback=litblue,coltext=white,colframe=litblue,arc=2mm]
    {\Large\bfseries Apunte #1 — #2}\\[3pt]
    {\small Proyecto Froth · aprender Machine Learning construyendo}
  \end{tcolorbox}\vspace{0.8em}}

% Cajas temáticas
\newtcolorbox{concepto}[1]{colback=blue!4,colframe=litblue,title={Concepto: #1},arc=1.5mm}
\newtcolorbox{practica}{colback=green!5,colframe=litgreen,title={Pruébalo tú (teclado en mano)},arc=1.5mm}
\newtcolorbox{ojo}{colback=orange!8,colframe=litorange,title={Ojo / error típico},arc=1.5mm}
\newtcolorbox{resumen}{colback=gray!8,colframe=gray!60,title={En una frase},arc=1.5mm}

\newcommand{\codigo}[1]{\texttt{#1}}
 justo después de \documentclass.
\usepackage[margin=2.4cm]{geometry}
\usepackage{xcolor}
\definecolor{litblue}{RGB}{31,79,121}
\definecolor{litgreen}{RGB}{27,94,32}
\definecolor{litorange}{RGB}{230,81,0}
\usepackage{parskip}
\usepackage{titlesec}
\titleformat{\section}{\Large\bfseries\color{litblue}}{\thesection}{0.6em}{}
\titleformat{\subsection}{\large\bfseries\color{litblue}}{\thesubsection}{0.6em}{}
\usepackage{tcolorbox}
\usepackage{fancyhdr}
\pagestyle{fancy}
\fancyhf{}
\fancyhead[L]{\small\color{litblue}Froth — Apuntes de ML}
\fancyhead[R]{\small\thepage}
\renewcommand{\headrulewidth}{0.4pt}
\usepackage[colorlinks=true,linkcolor=litblue,urlcolor=litblue]{hyperref}

% Cabecera de cada apunte: \cabecera{numero}{titulo}
\newcommand{\cabecera}[2]{%
  \begin{tcolorbox}[colback=litblue,coltext=white,colframe=litblue,arc=2mm]
    {\Large\bfseries Apunte #1 — #2}\\[3pt]
    {\small Proyecto Froth · aprender Machine Learning construyendo}
  \end{tcolorbox}\vspace{0.8em}}

% Cajas temáticas
\newtcolorbox{concepto}[1]{colback=blue!4,colframe=litblue,title={Concepto: #1},arc=1.5mm}
\newtcolorbox{practica}{colback=green!5,colframe=litgreen,title={Pruébalo tú (teclado en mano)},arc=1.5mm}
\newtcolorbox{ojo}{colback=orange!8,colframe=litorange,title={Ojo / error típico},arc=1.5mm}
\newtcolorbox{resumen}{colback=gray!8,colframe=gray!60,title={En una frase},arc=1.5mm}

\newcommand{\codigo}[1]{\texttt{#1}}

\begin{document}
\cabecera{16}{El bubble map interactivo (Plotly) y el truco del slider}

\section{Qué construimos}

El entregable estrella de la Fase 3: el \textbf{mapa de burbujas interactivo} de tus 126
papers. Existe en dos formas, alimentadas por el mismo código:
\begin{enumerate}
  \item \codigo{3\_Resultados/topic\_map.html} --- archivo autocontenido: doble clic y se abre
        en cualquier navegador, sin servidor. Compartible.
  \item La pestaña \textbf{Topic map} de tu web app (\codigo{app.py}), con controles en vivo.
\end{enumerate}

Qué se ve: un punto por paper, un color por subtema (geología azul, flotación verde,
economía/ajeno naranja, ruido gris), tamaño $=$ citas, la etiqueta de cada subtema flotando
sobre su nube, y al pasar el mouse: título, autores, año, citas.

\begin{concepto}{Plotly}
Librería de gráficos \textbf{interactivos}: en vez de una imagen fija (matplotlib), produce
HTML+JavaScript con zoom, pan, tooltips y leyenda clicable ``gratis''. Por eso el mapa se
puede guardar como archivo y abrirse en cualquier lado: el navegador hace el trabajo.
\end{concepto}

\begin{concepto}{una figura, dos caras (arquitectura)}
La función \codigo{bubble\_figure()} en \codigo{froth/visualize.py} construye la figura, y
DOS clientes la reutilizan: \codigo{plot\_bubbles()} la guarda como HTML, y \codigo{app.py}
la muestra en Streamlit. Si mañana mejoras el mapa, mejora en los dos sitios a la vez.
Principio general: \emph{el motor no se duplica; se envuelve}.
\end{concepto}

\section{El truco del slider de granularidad}

La web app tiene un slider ``min papers per subtopic'' que re-agrupa el mapa \textbf{al
instante}. ¿Cómo puede ser tan rápido, si el pipeline tarda?

\begin{concepto}{separar lo lento de lo rápido}
El pipeline tiene dos partes muy distintas:
\begin{itemize}
  \item \textbf{Lenta}: UMAP (acomodar 126 puntos desde 768D). Se calcula UNA vez y se cachea.
  \item \textbf{Rápida}: HDBSCAN sobre los puntos 2D ya acomodados (milisegundos) + re-etiquetar.
\end{itemize}
La clave: la granularidad \textbf{no afecta a UMAP}, solo a HDBSCAN. Así que mover el slider
solo repite la parte barata. Identificar qué se puede cachear y qué debe recalcularse es una
habilidad central de ingeniería de datos --- la acabas de usar.
\end{concepto}

\begin{ojo}
La granularidad es una \textbf{decisión de lectura}, no una verdad: con valor bajo verás
muchos subtemas pequeños (más detalle, más fragmentación); con valor alto, pocos grandes
(más síntesis). No existe ``el número correcto de clusters'' — existe el que mejor cuenta la
historia para tu propósito. Por eso es un slider tuyo y no una constante escondida.
\end{ojo}

\section{Dónde estamos frente a la competencia}

Con este mapa alcanzamos la \textbf{paridad} visual/interactiva (hover, zoom, filtros, tamaño
por citas) con Connected Papers, Litmaps o VOSviewer. Y ya tenemos diferenciadores que ellos
no: mapa por \emph{contenido} (no citas), subtemas auto-nombrados, granularidad en vivo,
ruido separable, búsqueda semántica integrada y todo open-source. Los diferenciadores
estrella (gaps + borrador de review + módulo tarea/PDFs) llegan en las Fases 4--5. La tabla
completa vive en \codigo{CLAUDE.md}.

\begin{resumen}
Plotly convierte el topic map en un mapa interactivo con dos caras (HTML compartible y
pestaña de la app) desde una sola función; el slider de granularidad es instantáneo porque
solo repite la parte barata (HDBSCAN 2D), no la cara (UMAP). Paridad con la competencia:
lograda; los diferenciadores grandes vienen en F4--F5.
\end{resumen}

\end{document}
